\section{Attention alternatives：从 full attention 到 hybrid state}

本节解释 linear attention、recurrent form、Mamba-2、Gated Delta Net、hybrid attention 和 DeepSeek Sparse Attention。


\begin{importantbox}{术语消化：linear attention、state、hybrid、DSA}
\begin{itemize}
    \item \textbf{Linear attention}：把 softmax attention 改写成可重排的核形式，使计算可用前缀状态累积。
    \item \textbf{Recurrent/state form}：维护一个随位置更新的 state，生成或长上下文时不必显式保存所有 pairwise attention。
    \item \textbf{Hybrid attention}：把 linear/state layers 和 full attention layers 交替使用，兼顾成本和质量。
    \item \textbf{DSA}：DeepSeek Sparse Attention，先筛选重要历史位置，再做稀疏 attention。
\end{itemize}
\end{importantbox}
\[
\mathrm{Attn}(Q,K,V)=\mathrm{softmax}(QK^\top)V.
\]
线性/稀疏变体都在改写这个式子的成本结构：要么避免完整 \(n\times n\) attention matrix，要么只选少量历史位置。

\begin{figure}[H]
\centering
\includegraphics[width=0.88\textwidth]{slide-003.jpg}
\caption{Slide 4: Linear attention}
\figsource{CS336 Spring 2026 Lecture 4 官方幻灯片。}
\end{figure}

\noindent\textbf{展开说明：}Consider the usual attention operation: $Q\in\mathbb{R}^{n\times d_k}$, $K\in\mathbb{R}^{n\times d_k}$, $V\in\mathbb{R}^{n\times d_v}$.

\begin{knowledgebox}{读图：Slide 4 应该怎么看}
这页是机制或证据页。先看它比较的是 compute、参数量、routing、负载均衡还是通信；再看结论支持的是质量提升、训练速度、推理成本还是系统复杂度。MoE 和 attention alternatives 的图常把模型结构和系统代价混在一起，读图时要分清“算法机制”和“硬件执行成本”。
\end{knowledgebox}


\begin{figure}[H]
\centering
\includegraphics[width=0.88\textwidth]{slide-004.jpg}
\caption{Slide 5: Recurrent form of linear attention}
\figsource{CS336 Spring 2026 Lecture 4 官方幻灯片。}
\end{figure}

\noindent\textbf{展开说明：}Recall that in purely linear attention, we consider the reordering

\begin{knowledgebox}{读图：Slide 5 应该怎么看}
这页是机制或证据页。先看它比较的是 compute、参数量、routing、负载均衡还是通信；再看结论支持的是质量提升、训练速度、推理成本还是系统复杂度。MoE 和 attention alternatives 的图常把模型结构和系统代价混在一起，读图时要分清“算法机制”和“硬件执行成本”。
\end{knowledgebox}


\begin{figure}[H]
\centering
\includegraphics[width=0.88\textwidth]{slide-005.jpg}
\caption{Slide 6: Minimax M1}
\figsource{CS336 Spring 2026 Lecture 4 官方幻灯片。}
\end{figure}

\noindent\textbf{展开说明：}Minimax M1 (and minimax-text-01) use a 7-to-1 hybrid (7 linear, 1 full) linear attention.


\begin{figure}[H]
\centering
\includegraphics[width=0.88\textwidth]{slide-006.jpg}
\caption{Slide 7: From linear attention to Mamba-2}
\figsource{CS336 Spring 2026 Lecture 4 官方幻灯片。}
\end{figure}

\noindent\textbf{展开说明：}Let’s generalize linear attention a little bit and add per-position weights..

\begin{knowledgebox}{读图：Slide 7 应该怎么看}
这页是机制或证据页。先看它比较的是 compute、参数量、routing、负载均衡还是通信；再看结论支持的是质量提升、训练速度、推理成本还是系统复杂度。MoE 和 attention alternatives 的图常把模型结构和系统代价混在一起，读图时要分清“算法机制”和“硬件执行成本”。
\end{knowledgebox}


\begin{figure}[H]
\centering
\includegraphics[width=0.88\textwidth]{slide-007.jpg}
\caption{Slide 8: Nemotron 3}
\figsource{CS336 Spring 2026 Lecture 4 官方幻灯片。}
\end{figure}

\noindent\textbf{展开说明：}Mamba attention hybrid (3-1 ish) – comparable (or better) pref to other similar models


\begin{figure}[H]
\centering
\includegraphics[width=0.88\textwidth]{slide-008.jpg}
\caption{Slide 9: Gated delta net (and friends)}
\figsource{CS336 Spring 2026 Lecture 4 官方幻灯片。}
\end{figure}

\noindent\textbf{展开说明：}Let’s generalize things further – gate the input and selectively erase the state.

\begin{knowledgebox}{读图：Slide 9 应该怎么看}
这页是机制或证据页。先看它比较的是 compute、参数量、routing、负载均衡还是通信；再看结论支持的是质量提升、训练速度、推理成本还是系统复杂度。MoE 和 attention alternatives 的图常把模型结构和系统代价混在一起，读图时要分清“算法机制”和“硬件执行成本”。
\end{knowledgebox}


\begin{figure}[H]
\centering
\includegraphics[width=0.88\textwidth]{slide-009.jpg}
\caption{Slide 10: Qwen 3.5 / Qwen Next}
\figsource{CS336 Spring 2026 Lecture 4 官方幻灯片。}
\end{figure}

\noindent\textbf{展开说明：}The newest qwen are 3-1 GDN / Attention hybrids.


\begin{figure}[H]
\centering
\includegraphics[width=0.88\textwidth]{slide-010.jpg}
\caption{Slide 11: Hybrid performance}
\figsource{CS336 Spring 2026 Lecture 4 官方幻灯片。}
\end{figure}

\noindent\textbf{展开说明：}Not many controlled ablations, but some evidence of low losses at small hybrid ratios

\begin{knowledgebox}{读图：Slide 11 应该怎么看}
这页是机制或证据页。先看它比较的是 compute、参数量、routing、负载均衡还是通信；再看结论支持的是质量提升、训练速度、推理成本还是系统复杂度。MoE 和 attention alternatives 的图常把模型结构和系统代价混在一起，读图时要分清“算法机制”和“硬件执行成本”。
\end{knowledgebox}


\begin{figure}[H]
\centering
\includegraphics[width=0.88\textwidth]{slide-011.jpg}
\caption{Slide 12: Alternative to hybrids: sparse adaptation}
\figsource{CS336 Spring 2026 Lecture 4 官方幻灯片。}
\end{figure}

\noindent\textbf{展开说明：}Instead of attending to every token, do sparse attention (DSA)


\begin{figure}[H]
\centering
\includegraphics[width=0.88\textwidth]{slide-012.jpg}
\caption{Slide 13: DSA – Deepseek Sparse Attention (v3.2, GLM5)}
\figsource{CS336 Spring 2026 Lecture 4 官方幻灯片。}
\end{figure}

\noindent\textbf{展开说明：}DSA – Deepseek Sparse Attention (v3.2, GLM5)

\begin{knowledgebox}{读图：Slide 13 应该怎么看}
这页是机制或证据页。先看它比较的是 compute、参数量、routing、负载均衡还是通信；再看结论支持的是质量提升、训练速度、推理成本还是系统复杂度。MoE 和 attention alternatives 的图常把模型结构和系统代价混在一起，读图时要分清“算法机制”和“硬件执行成本”。
\end{knowledgebox}

\subsection{本章小结}
本节的共同线索是 selective computation：通过结构选择让 token 不必总是访问所有历史或所有参数。但选择性越强，越需要额外机制保证信息流、负载均衡和训练稳定。

